% !TEX root = ../main.tex
\addchap*{Abstract}
\small
Protein structure prediction now spans entire proteomes, yet most AlphaFold models lack bound ligands and cofactors. This work builds and evaluates a reproducible pipeline that enriches AlphaFold monomers through DoGSite3 pocket detection, SIENA donor retrieval and alignment, ligand extraction, transplantation, and optional JAMDAscorer pose-only refinement. Three site-centric indicators are reported: local RMSD, TCS, and LEV, with a shared success gate of local RMSD $\leq$ \SI{1.0}{\angstrom} and TCS $\leq 0.35$ while LEV remains descriptive.

On a stratified benchmark of \num{2370} targets, \num{1840} yielded at least one placement, producing \num{6658} post-JAMDAscorer placements with a median of two per completed target. The joint gate retained \num{3592} placements (\SI{54.0}{\percent}) across \num{1457} targets, and LEV was defined for \num{6004} of \num{6658} placements overall and \num{5860} of \num{6488} organic placements. External benchmarking identified \num{2496} shared placements with AlphaFill; among the \num{1372} pairs where both Local RMSD and TCS were defined, pass rates were \SI{62.9}{\percent} (\num{863}/\num{1372}) for this pipeline and \SI{88.9}{\percent} (\num{1220}/\num{1372}, 95\%\,CI: 87.2--90.5) for AlphaFill, with moderate correlation in TCS and weaker alignment on local RMSD. Median wall times from instrumented runs were \SI{0.94}{\second} for DoGSite3, \SI{0.45}{\second} for SIENA, \SI{0.41}{\second} for JAMDAscorer, and \SI{6.0}{\second} per target overall, demonstrating scalable enrichment. Brief ablations indicate limited benefit from tightening the active-site identity threshold beyond 0.85, so the combined gate highlights residual errors from pocket-state mismatches, metal coordination differences, and low-identity outliers.

\normalsize
